\chapter{Perturbation theory}
\section{General methodology}
In quantum mechanics, we often solve problems by finding solutions to simpler systems and then determining how these solutions change when the system is modified. When the modification is small, we call it a perturbation.


The perturbation approach begins by expressing the Hamiltonian as:

\begin{equation}
H=H_0+\tau V
\end{equation}

Where $H_0$ is the unperturbed Hamiltonian with known eigenvalues and eigenstates:

\begin{equation}
H_0|n\rangle=E_n^0|n\rangle
\end{equation}

The perturbation is considered small when:

\begin{equation}
\langle n|H_0|n\rangle=E_n^0 \gg \tau\langle n|V|n\rangle
\end{equation}

Our goal is to solve the perturbed eigenvalue problem:

\begin{equation}
H|E_n\rangle=E_n|E_n\rangle
\end{equation}

We assume that both eigenstates and eigenvalues can be expanded in powers of $\tau$:

\[
\left\{\begin{array}{l}
|E_n\rangle=|a_n\rangle+\tau|b_n\rangle+\tau^2|c_n\rangle+\ldots \\
E_n=E_n^0+\tau E_n^1+\tau^2 E_n^2+\ldots
\end{array}\right.
\]

With $|a_n\rangle=|n\rangle$ so that when $\tau\rightarrow 0$ we recover the unperturbed solution. Substituting these expansions into the eigenvalue equation:

\[
\begin{array}{r}
H|E_n\rangle=(H_0+\tau V)(|n\rangle+\tau|b_n\rangle+\tau^2|c_n\rangle+\ldots)=\\
=H_0|n\rangle+\tau(H_0|b_n\rangle+V|n\rangle)+\tau^2(H_0|c_n\rangle+V|b_n\rangle)+\ldots
\end{array}
\]

For the right side of the eigenvalue equation:

\begin{gather}
E_n|E_n\rangle=(E_n^0+\tau E_n^1+\tau^2 E_n^2+\ldots)(|n\rangle+\tau|b_n\rangle+\tau^2|c_n\rangle+\ldots)=\\
=E_n^0|n\rangle+\tau(E_n^0|b_n\rangle+E_n^1|n\rangle)+\tau^2(E_n^0|c_n\rangle+E_n^1|b_n\rangle+E_n^2|n\rangle)+\ldots
\end{gather}

Equating coefficients of like powers of $\tau$:

\begin{align}
&H_0|n\rangle=E_n^0|n\rangle\\
&H_0|b_n\rangle+V|n\rangle=E_n^0|b_n\rangle+E_n^1|n\rangle \\
&H_0|c_n\rangle+V|b_n\rangle=E_n^0|c_n\rangle+E_n^1|b_n\rangle+E_n^2|n\rangle
\end{align}

We expand the correction states in terms of the unperturbed eigenbasis, excluding the state $|n\rangle$:

\begin{equation}
|b_n\rangle=\sum_{i\neq n}B_i(n)|i\rangle
\end{equation}

This exclusion is necessary because any component along $|n\rangle$ can be absorbed into the eigenvalue:

\begin{equation}
|b_n\rangle'=\sum_i B_i'(n)|i\rangle=B_n'(n)|n\rangle+\sum_{i\neq n}B_i(n)|i\rangle=B_n'(n)|n\rangle+|b_n\rangle
\end{equation}

This principle applies similarly to higher-order corrections. If we include such terms, the eigenvalue equation becomes:

\begin{align}
&H(|E_n\rangle+|n\rangle+\tau B_n'(n)|n\rangle+\tau^2 C_n'(n)|n\rangle+\ldots)=\\
&E_n(|E_n\rangle+|n\rangle+\tau B_n'(n)|n\rangle+\tau^2 C_n'(n)|n\rangle+\ldots)\\
&H|E_n\rangle+H(|n\rangle+\tau B_n'(n)|n\rangle+\tau^2 C_n'(n)|n\rangle+\ldots)=\\
&E_n|E_n\rangle+E_n(|n\rangle+\tau B_n'(n)|n\rangle+\tau^2 C_n'(n)|n\rangle+\ldots)
\end{align}

The action of $H$ on these additional terms gives:

\begin{align}
&H(|n\rangle+\tau B_n'(n)|n\rangle+\tau^2 C_n'(n)|n\rangle+\ldots)=\\
&(H|n\rangle+\tau B_n'(n)H|n\rangle+\tau^2 C_n'(n)H|n\rangle+\ldots)=\\
&(E_n|n\rangle+\tau B_n'(n)E_n|n\rangle+\tau^2 C_n'(n)E_n|n\rangle+\ldots)
\end{align}


These additional terms indeed cancel out.

For the first-order energy correction, we multiply the first-order equation by $\langle n|$:

\begin{equation}
\langle n|H_0|b_n\rangle+\langle n|V|n\rangle=\langle n|E_n^0|b_n\rangle+\langle n|E_n^1|n\rangle
\end{equation}

Since $|b_n\rangle$ is constructed to have no component along $|n\rangle$:

\begin{equation}
\langle n|b_n\rangle=\langle n|\sum_{i\neq n}B_i(n)|i\rangle=\sum_{i\neq n}B_i(n)\langle n|i\rangle=0
\end{equation}

The first-order correction simplifies to:

\begin{align}
&\langle n|H_0|b_n\rangle+\langle n|V|n\rangle=\langle n|E_n^0|b_n\rangle+\langle n|E_n^1|n\rangle\\
&\langle n|\sum_{i\neq n}B_i(n)H_0|i\rangle+\langle n|V|n\rangle=E_n^0\langle n|b_n\rangle+E_n^1\underbrace{\langle n|n\rangle}_{=1}\\
&\sum_{i\neq n}E_i^0B_i(n)\langle n|i\rangle+\langle n|V|n\rangle=E_n^1 \\
&\langle n|V|n\rangle=E_n^1
\end{align}

For the second-order energy correction:

\begin{align}
&\langle n|H_0|c_n\rangle+\langle n|V|b_n\rangle=\langle n|E_n^0|c_n\rangle+\langle n|E_n^1|b_n\rangle+\langle n|E_n^2|n\rangle\\
&\langle n|V|b_n\rangle=E_n^2
\end{align}

To find the corrections to the eigenstates, we project the first-order equation onto $\langle k|$ where $k\neq n$:

\begin{align}
&\langle k|H_0|b_n\rangle+\langle k|V|n\rangle=\langle k|E_n^0|b_n\rangle+\langle k|E_n^1|n\rangle\\
&\langle k|\sum_{i\neq n}B_i(n)H_0|i\rangle+\langle k|V|n\rangle=E_n^0\sum_{i\neq n}B_i(n)\underbrace{\langle k|i\rangle}_{\delta_{ik}}+E_n^1\langle k|n\rangle\\
&\sum_{i\neq n}E_i^0B_i(n)\underbrace{\langle k|i\rangle}_{\delta_{ik}}+\langle k|V|n\rangle=E_n^0B_k(n) \\
&(E_k^0-E_n^0)B_k(n)+\langle k|V|n\rangle=0\\
&B_k(n)=-\frac{\langle k|V|n\rangle}{E_k^0-E_n^0}
\end{align}

Similarly for higher-order corrections:

\begin{align}
C_k(n)&=-\frac{\langle k|V|b_n\rangle}{E_k^0-E_n^0}+E_n^1\frac{\langle k|b_n\rangle}{E_k^0-E_n^0} \\
D_k(n)&=-\frac{\langle k|V|c_n\rangle}{E_k^0-E_n^0}+E_n^1\frac{\langle k|c_n\rangle}{E_k^0-E_n^0}+E_n^2\frac{\langle k|b_n\rangle}{E_k^0-E_n^0}
\end{align}

\section{Application to the harmonic oscillator}
Consider a harmonic oscillator with Hamiltonian:

\begin{equation}
H=\frac{p^2}{2m}+\frac{k}{2}x^2
\end{equation}

If the spring constant undergoes a small change:

\begin{equation}
k\rightarrow k_0(1+\tau)\Rightarrow H=H_0+\tau\frac{k_0}{2}x^2
\end{equation}

The first-order energy correction is:

\begin{align}
E_n^1&=\langle n|V|n\rangle=\frac{k_0}{2}\langle n|x^2|n\rangle=\frac{k_0}{2}\frac{\hbar}{2m\omega_0}\langle n|(a^2+(a^\dagger)^2+2n+1)|n\rangle\\
&=\frac{m\omega_0^2}{2}\frac{\hbar}{m\omega_0}(n+\frac{1}{2})=\frac{\hbar\omega_0}{2}(n+\frac{1}{2})
\end{align}

For the second-order correction, we need:

\begin{align}
&|b_n\rangle=\sum_{k\neq n}B_k(n)|k\rangle \\
&B_k(n)=-\frac{\langle k|V|n\rangle}{E_k^0-E_n^0}
\end{align}


Substituting the expression for the first-order state correction:

\begin{equation}
|b_n\rangle=-\sum_{k\neq n}\frac{\langle k|V|n\rangle}{E_k^0-E_n^0}|k\rangle
\end{equation}

The second-order energy correction is:

\begin{equation}
E_n^2=\langle n|V|b_n\rangle
\end{equation}

Recognizing that $\langle k|V|n\rangle$ is generally complex and using $\langle n|V|k\rangle=\langle k|V|n\rangle^*$:

\begin{align}
E_n^2&=-\sum_{k\neq n}\langle n|V|\frac{\langle k|V|n\rangle}{E_k^0-E_n^0}k\rangle\\
&=-\sum_{k\neq n}\frac{\langle n|V|k\rangle\langle k|V|n\rangle}{E_k^0-E_n^0}\\
&=-\sum_{k\neq n}\frac{\langle k|V|n\rangle^*\langle k|V|n\rangle}{E_k^0-E_n^0} \\
&=-\sum_{k\neq n}\frac{|\langle k|V|n\rangle|^2}{E_k^0-E_n^0}
\end{align}

For the harmonic oscillator with perturbation $V=\frac{k_0}{2}x^2$, only states $n\pm 2$ contribute due to the selection rules of the ladder operators:

\begin{align}
E_n^2&=-\frac{|\langle n+2|V|n\rangle|^2}{E_{n+2}^0-E_n^0}-\frac{|\langle n-2|V|n\rangle|^2}{E_{n-2}^0-E_n^0}\\
&=-\frac{\hbar^2}{4m^2\omega_0^2}\frac{m^2\omega_0^4}{4}\frac{|\langle n+2|a^2|n\rangle|^2}{E_{n+2}^0-E_n^0}-\frac{\hbar^2}{4m^2\omega_0^2}\frac{m^2\omega_0^4}{4}\frac{|\langle n-2|(a^\dagger)^2|n\rangle|^2}{E_{n-2}^0-E_n^0}\\
&=-\frac{\hbar^2\omega_0^2}{16}\frac{(n+1)(n+2)|\langle n+2|n+2\rangle|^2}{E_{n+2}^0-E_n^0}-\frac{\hbar^2\omega_0^2}{16}\frac{n(n-1)|\langle n-2|n-2\rangle|^2}{E_{n-2}^0-E_n^0}\\
&=-\frac{\hbar^2\omega_0^2}{16}\frac{(n+1)(n+2)}{2\hbar\omega_0}-\frac{\hbar^2\omega_0^2}{16}\frac{n(n-1)}{-2\hbar\omega_0}\\
&=-\frac{\hbar\omega_0}{32}(n^2+3n+2-n^2+n)=-\frac{\hbar\omega_0}{8}(n+\frac{1}{2})
\end{align}

Continuing this process for higher-order corrections, we obtain:

\begin{equation}
E_n=\hbar\omega_0(n+\frac{1}{2})(1+\frac{\tau}{2}-\frac{\tau^2}{8}+\ldots)
\end{equation}

This series is precisely the Taylor expansion of:

\begin{equation}
\sqrt{1+\tau}\sim 1+\frac{\tau}{2}-\frac{\tau^2}{8}+\ldots
\end{equation}

The exact result should be $\hbar\omega=\hbar\frac{\sqrt{k}}{m}$, and with the perturbed spring constant:

\begin{equation}
\hbar\omega=\hbar\frac{\sqrt{k_0(1+\tau)}}{m}=\hbar\omega_0\sqrt{1+\tau}
\end{equation}

This confirms that our perturbation approach correctly reproduces the exact solution.
